\scp{Eastern Indonesia and Timor-Leste as a zone of contact}{01-06}
In the previous section it was noted that early scholars observed an affinity
between the Austronesian languages of our target region with those of “Melanesia”,
and some even referred to it as “the beginning of Melanesia”.
There is an increasing awareness that many features that set these languages apart from those to the west
and north are due to contact with lexically
and typologically different SOV Papuan languages still found in the region.
And a growing awareness that prior to the arrival of Austronesian-speaking peoples,
there is evidence that Papuan languages were much
more widespread in the region than they are today
\citep{Grimes1991a,Grimes2000a,GrimesCIP,GrimesTherikGrimesJacob1997,
KlamerReesinkVanStaden2008,DonohueGrimes2008,Schapper2015,Schapper2024,Klamer2019}.

As we show in this study,
the complex and diverse on-the-ground data we find in the region suggest
that speakers of multiple Austronesian languages
(probably more diverse than a single linkage)
came into contact with speakers of diverse Papuan languages already in place.
There arose intense contact, bilingualism
and intermarriage over millennia resulting in structural
and typological accommodation in the languages.
The further development of Austronesian languages
and patterns of contact were not uniform throughout the region.
Rather than a single monolithic Austronesian incursion,
there were very likely multiple incursions (perhaps during different time periods
and from different starting places),
many internal migrations,
and trade spanning more than two millennia (see \srf{02-02-01}),
which all add further complexities to unravelling the prehistory of the region.

\citet{Schapper2015} has coined the term ``Linguistic Wallacea''
for this region of eastern Indonesia
and Timor-Leste that includes non-Oceanic Austronesian languages
showing varying degrees of contact-induced change from Papuan languages.\footnote{
		\citet[59]{Oppenheimer2003} refers to “the languages of Wallacea”
		as sandwiched between Oceanic languages to the east and others to the west.
		He notes that the genetics of the Wallacean region figure significantly
		in unravelling the prehistory of the spread of Austronesian-speaking populations.}
See \frf{01-05}.
Linguistic Wallacea is the westernmost extent of the area identified by \citet{Schapper2015}
as Linguistic Melanesia, and in a very real sense this
zone of contact is the beginning of Melanesia when travelling from west to east.
The focus of this study is on the Austronesian languages of “Linguistic Wallacea”
(excluding the Eastern Malayo-Polynesian languages of south Halmahera and Indonesian Papua).
The subgroups within this region will be referred to as “Wallacean” subgroups,
indicating a geographical region,
and not making any claims about linguistic relatedness at an intermediate level.

\begin{figure}[h]
	\includegraphics[width=\columnwidth]{Wallacea.pdf}
	\fcp{Linguistic Wallacea}{01-05}
\end{figure}

The many Papuan-like features found in Austronesian languages in eastern Indonesia
and Timor-Leste (\srf{02-03-01}) are actually much more geographically widespread
than the limited locations in which Papuan languages are still spoken today (see \frf{01-05}).
This can be seen as trace evidence that when the Austronesian-speaking peoples arrived in the region,
they encountered Papuan-speaking populations throughout the region.\footnote{
		The existence into historical times of the non-Austronesian language Tambora
		on Sumbawa Island \citep{Donohue2007} provides direct evidence that
		other Papuan languages
		were once much more widespread in the region than the geographically
		more limited \tsc{Timor-Alor-Pantar}, \tsc{North Halmahera}
		and other Papuan languages in the Bird’s Neck of Papua still present.}

Hereafter, unless otherwise specified, our use of the term “the region”,
“this region” or “our target region” refers to the parts of present
day eastern Indonesia and Timor-Leste between the Blust line and EMP.
See \frf{01-06}.

\begin{figure}[h]
	\includegraphics[width=\columnwidth]{families.pdf}
	\fcp{Eastern Indonesia and Timor-Leste with subgroups}{01-06}
\end{figure}